\documentclass[10pt,a4paper]{amsart}
\usepackage[margin=19mm]{geometry}
\usepackage{amsmath,amssymb,mathtools}
\usepackage[scr=rsfs,cal=euler]{mathalfa}
\usepackage{xfrac}
\usepackage[unicode,hidelinks,bookmarks=false]{hyperref}
\newcommand{\ie}{i.e.\ }
\newcommand{\cf}{cf.\ }
\newcommand{\resp}{respectively\ }
\newcommand{\Subsection}{\subsection}
\providecommand{\nobreakdash}{\nobreak\hbox{-}}
\setlength{\emergencystretch}{2em}
\multlinegap0pt
\newcommand{\E}{\mathbb E}
\newcommand{\Id}{I}
\newcommand{\KL}{\operatorname{KL}}
\newcommand{\R}{\mathbb R}
\newcommand{\cH}{\mathcal H}
\newcommand{\cI}{\mathcal I}
\newcommand{\dd}{\mathrm d}
\newcommand{\eps}{\varepsilon}
\theoremstyle{plain}
\newtheorem{theorem}{Theorem}[section]
\newtheorem{proposition}[theorem]{Proposition}
\newtheorem{lemma}[theorem]{Lemma}
\newtheorem{corollary}[theorem]{Corollary}
\theoremstyle{definition}
\newtheorem{assumption}[theorem]{Assumption}
\newtheorem{remark}[theorem]{Remark}
\title[Global Convergence of Third-Order Langevin Dynamics for Non-Convex Optimization via Simulated Annealing]{Global Convergence of Third-Order Langevin Dynamics for Non-Convex Optimization via Simulated Annealing}
\author{Yingli Wang, Lingjiong Zhu}
\date{Source: arXiv:2609.28611v3 (CC BY 4.0)}
\begin{document}
\maketitle

\noindent\emph{Source and licence.} The delivered work is the article shown in the title above, version arXiv:2609.28611v3 (CC BY 4.0), distributed under the Creative Commons Attribution 4.0 International licence, as recorded in the arXiv licence metadata for this exact version and shown on the abstract page saved with this item. The full licence notice, which carries the licence URL and the address of that abstract page, is the bundled file \texttt{licenses/arxiv-2609.28611-CC-BY-4.0.txt}.
\par\smallskip

\noindent\textit{Editorial scope.} Counted unit: the article's Theorem (source label thm:main) (source0.tex lines 1116--1126, source label \texttt{thm:main}): ``Continuous-time annealing Under Assumptions , , and , for every and there is such that, for all sufficiently large , P (U(X t) ) C , (t t 0) - E, ,(1-D E) 2 . Here is the cooling limit in , and is the''. It is delivered together with the article's complete original proof (source0.tex lines 1143--1442, one proof environment printed later in the article, detached from the statement by the article's own arrangement, 300 lines), copied line for line from the article's own text. Required context retained uncounted: eq:model (lines 392--399); eq:generator (lines 409--412); ass:potential (lines 432--446); ass:initial (lines 451--461); ass:lsi (lines 463--477); eq:tensorizedlsi (lines 493--495); eq:matrix (lines 565--569); lem:matrix (lines 597--608); lem:closure (lines 650--679). Editorial formatting only: an AMS \texttt{amsart} preamble that restates the article's own macro definitions and its own theorem declarations, and the theorem environments the retained lines use; the article's own class file is neither shipped nor input; the retained environments are renumbered sequentially in order of appearance, whereas the article prints them with its own numbering; the article's equation numbering is not reproduced and every cross-reference inside the retained text resolves to the wrapper's numbering. No citation occurs inside any retained line, so the wrapper carries no bibliography; All retained bodies are exact copies of the bundled \texttt{sources/211603/source0.tex}, so no omission receipt applies. No asset-inclusion command occurs inside any retained span and the package ships no image asset.


\section*{Retained context: eq:model (source0.tex lines 392--399)}

\begin{equation}\label{eq:model}
\begin{aligned}
 \dd X_t&=V_t\dd t,\\
 \dd V_t&=(-\nabla U(X_t)+\lambda Z_t)\dd t,\\
 \dd Z_t&=(-\lambda V_t-\gamma Z_t)\dd t
                 +\sqrt{2\gamma\eps_t}\,\dd B_t,
\end{aligned}
\end{equation}


\section*{Retained context: eq:generator (source0.tex lines 409--412)}

\begin{equation}\label{eq:generator}
 L_\eps:=v\cdot\nabla_x+(-\nabla U+\lambda z)\cdot\nabla_v
       +(-\lambda v-\gamma z)\cdot\nabla_z+\gamma\eps\Delta_z.
\end{equation}


\section*{Retained context: ass:potential (source0.tex lines 432--446)}

\begin{assumption}[Potential and cooling]\label{ass:potential}
The function $U\in C^4(\R^d)$ has minimum zero, bounded Hessian
$\|\nabla^2U\|_{\rm op}\leq K$, and satisfies
$x\cdot\nabla U(x)\geq r|x|^2-m$ for some $r>0,m\geq0$.
Its third and fourth derivatives have at most polynomial growth.
The temperature $\eps_t\in C^1([0,\infty);(0,\infty))$ is non-increasing in $t$. For
some $t_0>1$ and $E>0$,
\begin{equation}\label{eq:coolingassumption}
 \eps_t\log(t+t_0)\longrightarrow E,
 \qquad
 |\eps_t'|(1+\eps_t^{-3})
 \leq C\frac{\log(t+t_0)}{t+t_0}
 \quad\text{for all sufficiently large }t.
\end{equation}
\end{assumption}


\section*{Retained context: ass:initial (source0.tex lines 451--461)}

\begin{assumption}[Continuous-time initial law]\label{ass:initial}
The initial law $\nu_0$ of \eqref{eq:model} is absolutely continuous, has a finite
second moment, and finite
relative Fisher information with respect to $\pi_{\eps_0}$:
\begin{equation}\label{eq:initialfisher}
 \cI(\nu_0\mid\pi_{\eps_0})
 :=\int_{\R^{3d}}
 \left|\nabla\log\frac{\dd\nu_0}{\dd\pi_{\eps_0}}\right|^2\dd\nu_0
 <\infty.
\end{equation}
\end{assumption}


\section*{Retained context: ass:lsi (source0.tex lines 463--477)}

\begin{assumption}[LSI barrier]\label{ass:lsi}
Let $Z_\eps:=\int_{\R^d}e^{-U(y)/\eps}\dd y$ and
$\mu_\eps(\dd x)=Z_\eps^{-1}e^{-U(x)/\eps}\dd x$ be the position
Gibbs measure. There is $D\geq0$ such that for every $\eta>0$ there are
$C_\eta<\infty$ and $\eps_\eta>0$ for which
\begin{equation}\label{eq:lsi}
 \KL(\nu\mid\mu_\eps)\leq C_{\rm LS}^{\rm pos}(\eps)
       \cI(\nu\mid\mu_\eps),
 \qquad C_{\rm LS}^{\rm pos}(\eps)\leq C_\eta e^{(D+\eta)/\eps},
 \quad 0<\eps\leq\eps_\eta.
\end{equation}
Here $\cI(\nu\mid\mu):=\int_{\R^d}|\nabla\log(\dd\nu/\dd\mu)|^2\dd\nu$, and we assume
$D<E$ where $E$ is the cooling
constant from Assumption~\ref{ass:potential}.
\end{assumption}


\section*{Retained context: eq:tensorizedlsi (source0.tex lines 493--495)}

\begin{equation}\label{eq:tensorizedlsi}
 C_{\rm LS}(\eps)=\max\{C_{\rm LS}^{\rm pos}(\eps),\eps/2\},
\end{equation}


\section*{Retained context: eq:matrix (source0.tex lines 565--569)}

\begin{align}\label{eq:matrix}
A_0:=\frac{2\lambda b+2(q_1^2+q_2^2)+1}{\gamma},
\qquad
M:=\begin{pmatrix}1&1&0\\1&2&b\\0&b&\varsigma\end{pmatrix}\otimes\Id_d,
\end{align}


\section*{Retained context: lem:matrix (source0.tex lines 597--608)}

\begin{lemma}[Algebraic dissipation certificate]\label{lem:matrix}
The matrix $M$ in \eqref{eq:matrix} is positive definite. For every symmetric $Q$ with
$\|Q\|_{\rm op}\leq K$, define
\[
 J_Q:=\begin{pmatrix}0&-\Id_{d}&0\\Q&0&-\lambda\Id_{d}\\0&\lambda\Id_{d}&-\gamma\Id_{d}\end{pmatrix},
 \qquad P_z:=\operatorname{diag}(0,0,\Id_d).
\]
Then
\begin{equation}\label{eq:certificate}
 -J_QM-MJ_Q^\top+A_0\gamma P_z\succeq\frac12\Id_{3d}.
\end{equation}
\end{lemma}


\section*{Retained context: lem:closure (source0.tex lines 650--679)}

\begin{lemma}[Justification by initial-density approximation]\label{lem:closure}
Under Assumption~\ref{ass:potential}, let the annealed third-order process
$S_t=(X_t,V_t,Z_t)$ in \eqref{eq:model} be started, or restarted
at a deterministic time $s\geq0$, from
an absolutely continuous law $\nu_s$ with finite relative Fisher information with
respect to $\pi_{\eps_s}$ and finite second moment. Let $\nu_t$ denote the
law of $S_t$ for $t\geq s$, and let $F(t):=\cH_{\eps_t}(\nu_t)$. Then
$F(t)<\infty$ for every finite $t\geq s$.
Moreover, there exist smooth, strictly positive probability laws
$\nu_s^{(n)}$, Gaussian outside compact sets, such that
\[
 \mathcal W_2\left(\nu_s^{(n)},\nu_s\right)\to0,\qquad
 \cH_{\eps_s}\left(\nu_s^{(n)}\right)\to F(s),\qquad
 \sup_{n\geq1}\int_{\R^{3d}}|y|^2\nu_s^{(n)}(\dd y)<\infty.
\]
Let $\nu_u^{(n)}$ be the laws of \eqref{eq:model} started from $\nu_s^{(n)}$
at time $s$, and set $F_n(u):=\cH_{\eps_u}\left(\nu_u^{(n)}\right)$.
Fix $T>s$.
Suppose that there exist $a,b\in L^1([s,T])$, with $b\geq0$, such that
\begin{equation}\label{eq:abstractcomparison}
 F_n'(u)+a(u)F_n(u)\leq b(u)
 \quad\text{for a.e. }u\in[s,T]\text{ and every }n\geq1.
\end{equation}
Then $F$ satisfies
\begin{equation}\label{eq:abstractintegrated}
 F(t)\leq e^{-\int_s^t a(u)\dd u}F(s)
 +\int_s^t e^{-\int_v^t a(u)\dd u}b(v)\dd v,
 \qquad s\leq t\leq T.
\end{equation}
\end{lemma}


\section{The counted unit: Theorem (source label thm:main) and its complete original proof}

\begin{theorem}[Continuous-time annealing]\label{thm:main}
Under Assumptions~\ref{ass:potential}, \ref{ass:initial}, and \ref{ass:lsi}, for every
$\delta>0$ and $\alpha>0$ there is $C_{\delta,\alpha}<\infty$ such that, for
all sufficiently large $t$,
\begin{equation}\label{eq:mainrate}
 \mathbb P\left(U(X_t)>\delta\right)
 \leq C_{\delta,\alpha}(t+t_0)^{-\min\{\delta/E,\,(1-D/E)/2\}+\alpha}.
\end{equation}
Here $E$ is the cooling limit in \eqref{eq:coolingassumption}, and $D$ is the
LSI barrier in \eqref{eq:lsi}.
\end{theorem}

\begin{proof}[Proof of Theorem~\ref{thm:main}]
We separate moment control, frozen dissipation, and the cooling contribution.

\emph{Moment control.}
For $a>0$ small, let $b_a=2a/\lambda$, $c_a=a\lambda/\gamma$, and
\[
 V_a=1+H+a x\cdot v+b_a v\cdot z+c_a x\cdot z
                       +\frac12c_a\lambda|x|^2+C_a.
\]
Choose the constant $C_a$ so that $V_a\geq1$. Dissipativity and bounded Hessian
give quadratic lower and upper bounds on $U$.
To make them explicit, set
$R_0:=\max\{1,\sqrt{2m/r}\}$. If $x=\rho\theta$ with $|\theta|=1$ and
$\rho\geq R_0$, then
\[
 \frac{\dd}{\dd\rho}U(\rho\theta)
 =\theta\cdot\nabla U(\rho\theta)
 \geq r\rho-\frac m\rho\geq\frac r2\rho.
\]
After integration from $R_0$ to $\rho$ and adjustment on the ball $B_{R_0}$,
this gives $U(x)\geq r|x|^2/4-C_-$. On the other hand, Taylor's formula and
$\|\nabla^2U\|_{\rm op}\leq K$ give
\[
 U(x)\leq U(0)+|\nabla U(0)||x|+\frac K2|x|^2
 \leq C_+(1+|x|^2).
\]
Thus,
\begin{equation}\label{eq:quadraticU}
 \frac r4|x|^2-C_-\leq U(x)\leq C_+(1+|x|^2).
\end{equation}
Consequently, for small $a$,
$V_a$ is comparable to $1+|x|^2+|v|^2+|z|^2$.
Direct calculation, with cancellation of the $x\cdot v$ and $x\cdot z$ terms,
gives
\begin{align*}
 L_\eps V_a={}&-a x\cdot\nabla U-a|v|^2+(2a-\gamma)|z|^2
 +(c_a-\gamma b_a)v\cdot z-b_a\nabla U\cdot z+\gamma\eps d.
\end{align*}
For clarity, let $c_0=|\nabla U(0)|$ and
$c_1=|\lambda/\gamma-2\gamma/\lambda|$. Young's inequality gives
\begin{align}\label{eq:lyapunovyoung}
 a c_1|v||z|&\leq\frac a4|v|^2+a c_1^2|z|^2,\\
 \frac{2a}{\lambda}(K|x|+c_0)|z|
 &\leq\frac{ar}{4}|x|^2
       +a\left(\frac{4K^2}{r\lambda^2}+1\right)|z|^2
       +\frac{a c_0^2}{\lambda^2}.
\end{align}
In addition to the condition ensuring quadratic comparability of $V_a$, take
$a(3+c_1^2+4K^2/(r\lambda^2))\leq\gamma/2$.
Using $x\cdot\nabla U(x)\geq r|x|^2-m$ together with
\eqref{eq:lyapunovyoung} gives
\[
 L_{\eps_t}V_a
 \leq-\frac{3ar}{4}|x|^2-\frac{3a}{4}|v|^2-\frac{\gamma}{2}|z|^2
       +am+\frac{ac_0^2}{\lambda^2}+\gamma d\sup_{u\geq0}\eps_u
 \leq-cV_a+C.
\]
Dynkin's formula and Gr\"{o}nwall's inequality now yield
\begin{equation}\label{eq:moments}
 \sup_{t\geq0}\E\left[|X_t|^2+|V_t|^2+|Z_t|^2\right]<\infty.
\end{equation}
Moreover $|\nabla_z V_a|^2\leq C V_a$, and hence
\begin{align*}
 L_{\eps_t}V_a^p
 =pV_a^{p-1}L_{\eps_t}V_a
   +\gamma\eps_t p(p-1)V_a^{p-2}|\nabla_zV_a|^2\leq-c_pV_a^p+C_p,
 \qquad p\in\mathbb N.
\end{align*}
In particular, $p=1$ proves the required second-moment bound. Higher moments
are propagated whenever the corresponding initial moment is finite. The
second-moment estimate is uniform for restarted processes with uniformly
bounded initial second moments.
For the Gibbs position marginal, integration by parts gives
$\pi_\eps(x\cdot\nabla U)=d\eps$. Dissipativity therefore bounds its
second moment uniformly for $0<\eps\leq\eps_0$. The quadratic upper bound
on $U$ and the two Gaussian factors then bound $\pi_\eps(H)$ uniformly.
Indeed,
\[
 0=\int_{\R^d}\nabla\cdot\left(xe^{-U(x)/\eps}\right)\dd x
 =dZ_\eps-\eps^{-1}\int_{\R^d}
 x\cdot\nabla U(x)e^{-U(x)/\eps}\dd x,
\]
and hence
\[
 r\mu_\eps(|x|^2)-m\leq d\eps,
 \qquad
 \mu_\eps(|x|^2)\leq\frac{d\eps+m}{r}.
\]
Together with \eqref{eq:quadraticU} and
$\pi_\eps(|v|^2)=\pi_\eps(|z|^2)=d\eps$, this yields
\begin{equation}\label{eq:gibbsmoments}
 \sup_{0<\eps\leq\eps_0}
 \left\{\pi_\eps\left(|x|^2+|v|^2+|z|^2\right)+\pi_\eps(H)\right\}<\infty.
\end{equation}

\emph{Frozen dissipation.}
Recall from \eqref{eq:generator} that
\[
 L_\eps=v\cdot\nabla_x+(-\nabla U+\lambda z)\cdot\nabla_v
 +(-\lambda v-\gamma z)\cdot\nabla_z+\gamma\eps\Delta_z.
\]
Its adjoint in $L^2(\pi_\eps)$
has drift
$b^\dagger=(-v,\nabla U-\lambda z,\lambda v-\gamma z)$, with Jacobian
$J_{\nabla^2U}$. Relative-density differentiation is governed by this adjoint
drift. Let $\nu_t^\eps$ denote the
law evolving with temperature fixed at $\eps$, and let
$h:=\dd\nu_t^\eps/\dd\pi_\eps$, $g:=\log h$, $w:=\nabla g$, and
$\mathcal S:=\gamma\eps\Delta_z-\gamma z\cdot\nabla_z$. Then
\[
 \partial_t g=L_\eps^\dagger g+\gamma\eps|\nabla_z g|^2,
 \qquad L_\eps+L_\eps^\dagger=2\mathcal S.
\]
Differentiate $\int_{\R^{3d}} h w^\top M w\dd\pi_\eps$, move $L_\eps^\dagger$
off $h$, and use
$\nabla(L_\eps^\dagger g)=L_\eps^\dagger w+J_{\nabla^2U}^\top w$.
Integration by parts with $\mathcal S$ cancels the term containing
$\nabla|\nabla_z g|^2$, leaving
\begin{subequations}\label{eq:fisherentropyderivatives}
\begin{align}
 \frac{\dd}{\dd t}\cI_M(\nu_t^\eps\mid\pi_\eps)
 &=\int_{\R^{3d}} w^\top\left(J_{\nabla^2U}M+MJ_{\nabla^2U}^\top\right)w\dd\nu_t^\eps
   -2\gamma\eps\sum_{i=1}^{d}\int_{\R^{3d}}(\partial_{z_i}w)^\top M(\partial_{z_i}w)\dd\nu_t^\eps,
   \label{eq:fisherderivative}\\
 \frac{\dd}{\dd t}\KL(\nu_t^\eps\mid\pi_\eps)
 &=-\gamma\eps\int_{\R^{3d}}|w_z|^2\dd\nu_t^\eps.
   \label{eq:entropyderivative}
\end{align}
\end{subequations}
These operations are justified first for the regular initial laws in
Lemma~\ref{lem:closure}. Multiplying the entropy identity
\eqref{eq:entropyderivative}
by $A_0/\eps$ and adding it to the weighted Fisher-information identity
\eqref{eq:fisherderivative}
gives
\begin{align}\label{eq:frozenidentity}
 \frac{\dd}{\dd t}\cH_\eps(\nu_t^\eps)
 =&-\int_{\R^{3d}} w^\top\left(-J_{\nabla^2U}M-MJ_{\nabla^2U}^\top
                         +A_0\gamma P_z\right)w\,\dd\nu_t^\eps\notag\\
 &\qquad\qquad-2\gamma\eps\sum_{i=1}^d\int_{\R^{3d}}
           (\partial_{z_i}w)^\top M(\partial_{z_i}w)\,\dd\nu_t^\eps.
\end{align}
Lemma~\ref{lem:matrix} implies
\[
 \frac{\dd}{\dd t}\cH_\eps(\nu_t^\eps)
 \leq-\frac12\int_{\R^{3d}}|w|^2\dd\nu_t^\eps.
\]
On the other hand, \eqref{eq:lsi} and tensorization as in
\eqref{eq:tensorizedlsi} give
\[
 \cH_\eps(\nu_t^\eps)
 \leq\left(\|M\|+\frac{A_0C_{\rm LS}(\eps)}{\eps}\right)
       \int_{\R^{3d}}|w|^2\dd\nu_t^\eps.
\]
Consequently,
\begin{equation}\label{eq:decay}
 \frac{\dd}{\dd t}\cH_\eps(\nu_t^\eps)
 \leq-\kappa(\eps)\cH_\eps(\nu_t^\eps),\qquad
 \kappa(\eps):=\frac{1}{2(\|M\|+A_0C_{\rm LS}(\eps)/\eps)}.
\end{equation}
To record the loss in this estimate explicitly, fix $\eta>0$ and first choose
$\eta_{\rm LSI}>0$ sufficiently small such that
$(D+\eta_{\rm LSI})/E<D/E+\eta/3$. Equations
\eqref{eq:lsi} and \eqref{eq:decay} give, for small $\eps$,
\begin{align}\label{exp:as:lower:bound}
 \kappa(\eps)\geq c\eps e^{-(D+\eta_{\rm LSI})/\eps}.
\end{align}
Since $\eps_t\log(t+t_0)\to E$, for all sufficiently large $t$,
\[
 \frac{D+\eta_{\rm LSI}}{\eps_t}
 \leq\left(\frac DE+\frac{2\eta}{3}\right)\log(t+t_0).
\]
Consequently, the exponential in \eqref{exp:as:lower:bound} satisfies
\[
 e^{-(D+\eta_{\rm LSI})/\eps_t}
 \geq(t+t_0)^{-D/E-2\eta/3}.
\]
Finally,
$\eps_t\asymp1/\log(t+t_0)\geq c_\eta(t+t_0)^{-\eta/3}$, and hence
\begin{equation}\label{eq:kappalower}
 \kappa(\eps_t)\geq c_\eta(t+t_0)^{-D/E-\eta}.
\end{equation}

\emph{Cooling contribution.}
For a fixed law $\nu$ with density $\rho$, write
$w:=\nabla\log(\rho/\pi_\eps)$. Then
$\partial_\eps\log(\rho/\pi_\eps)=-(H-\pi_\eps(H))/\eps^2$ and
$\partial_\eps w=-\nabla H/\eps^2$. Therefore,
\begin{align}\label{eq:temperatureidentity}
 \partial_\eps\cH_\eps(\nu)
 ={}&-\frac{2}{\eps^2}\int_{\R^{3d}} w^\top M\nabla H\,\dd\nu
 -\frac{A_0}{\eps^2}\KL(\nu\mid\pi_\eps)
 -\frac{A_0}{\eps^3}\left(\nu(H)-\pi_\eps(H)\right).
\end{align}
For the actual annealing law $\nu_t$, apply \eqref{eq:frozenidentity}
instantaneously at $\eps=\eps_t$ and add $\eps_t'$ times
\eqref{eq:temperatureidentity}. Set $F(t)=\cH_{\eps_t}(\nu_t)$. The moment
bounds \eqref{eq:moments} for $\nu_t$ and \eqref{eq:gibbsmoments} for
$\pi_{\eps_t}$, together with $|\nabla H(y)|\leq C(1+|y|)$ and the
Cauchy--Schwarz inequality, give
\[
 \left|\int_{\R^{3d}} w^\top M\nabla H\dd\nu_t\right|\leq C\sqrt{F(t)},
 \quad \KL(\nu_t\mid\pi_{\eps_t})\leq\eps_t F(t)/A_0,
 \quad |\nu_t(H)-\pi_{\eps_t}(H)|\leq C.
\]
Since $\sqrt F\leq1+F$, the second condition in
\eqref{eq:coolingassumption} bounds the cooling term by
\begin{equation}\label{eq:coolingexplicit}
 C|\eps_t'|\left(1+\eps_t^{-3}\right)(1+F(t))
 \leq C\frac{\log(t+t_0)}{t+t_0}(1+F(t)),
\end{equation}
for all sufficiently large $t$. The constant depends on the initial law
only through its second-moment bound. On every fixed finite interval the
first coefficient in \eqref{eq:coolingexplicit} is bounded. Dropping frozen
dissipation and applying Gr\"{o}nwall's inequality therefore also bounds $F$ before the
low-temperature LSI becomes available, uniformly over the approximations
in Lemma~\ref{lem:closure}.
More explicitly, for each $T<\infty$ there is $C_T<\infty$, depending only
on the common second-moment bound and the model parameters, such that, with
$F_n(t):=\cH_{\eps_t}\left(\nu_t^{(n)}\right)$,
\[
 F_n'(t)\leq C_T(1+F_n(t)),\qquad 0\leq t\leq T.
\]
Thus, for $0\leq s\leq t\leq T$,
\[
 F_n(t)
 \leq e^{C_T(t-s)}\left(F_n(s)+C_T(t-s)\right).
\]
The same $C_T$ applies to every $n$.

For any $0<\eta<1-D/E$, the ratio of
$\log(t+t_0)/(t+t_0)$ to $(t+t_0)^{-D/E-\eta}$ tends to zero.
Absorb the coefficient of $F$ in half of the frozen decay and use
$\log(t+t_0)\leq C_\eta(t+t_0)^\eta$. This gives
\begin{equation}\label{eq:ode}
 F'(t)\leq-c_\eta(t+t_0)^{-D/E-\eta}F(t)
              +C_\eta(t+t_0)^{-1+\eta},
 \qquad F(t)=\cH_{\eps_t}(\nu_t).
\end{equation}
For any $r_0<1-D/E$, choose $\eta$ so that
$r_0<1-D/E-2\eta$. A sufficiently large multiple of
$(t+t_0)^{-r_0}$ is a supersolution of \eqref{eq:ode}. Indeed, if
$\overline F(t)=C(t+t_0)^{-r_0}$, then
\begin{equation}\label{eq:supersolutioncalculation}
 \overline F'(t)+c_\eta(t+t_0)^{-D/E-\eta}\overline F(t)
 ={C}\left[-r_0(t+t_0)^{-r_0-1}
 +c_\eta(t+t_0)^{-r_0-D/E-\eta}\right].
\end{equation}
Because $r_0+D/E+\eta<1-\eta$ and $D/E+\eta<1$, the right-hand side is at
least $C_\eta(t+t_0)^{-1+\eta}$ for all large $t$ after increasing $C$.
Choose $t_1$ sufficiently large and increase $C$ once more so that
$F(t_1)\leq\overline F(t_1)$. For $G=F-\overline F$, subtracting the
supersolution inequality obtained from
\eqref{eq:supersolutioncalculation} from \eqref{eq:ode}
gives
\[
 G'(t)+c_\eta(t+t_0)^{-D/E-\eta}G(t)\leq0,
 \qquad t\geq t_1.
\]
Multiplication by the corresponding integrating factor yields
\[
 G(t)\leq G(t_1)
 \exp\left(-\int_{t_1}^tc_\eta(u+t_0)^{-D/E-\eta}\dd u\right)\leq0.
\]
Therefore
\[
F(t)\leq C_{r_0}(t+t_0)^{-r_0}.
\]

Finally, $\KL(\nu_t\mid\pi_{\eps_t})\leq\eps_tF(t)/A_0$.
Pinsker's inequality gives
\[
 \mathbb P(U(X_t)>\delta)
 \leq\pi_{\eps_t}(U>\delta)
                   +\sqrt{\frac12\KL(\nu_t\mid\pi_{\eps_t})}.
\]
The Gibbs tail estimate used here is
\begin{equation}\label{eq:gibbstail}
 \pi_\eps(U>\delta)\leq C_{\delta,\eta}e^{-(\delta-\eta)/\eps},
 \qquad 0<\eta<\delta.
\end{equation}
To verify \eqref{eq:gibbstail}, choose a measurable set $A_\eta$ of positive
Lebesgue measure on which $U\leq\eta/2$. For $0<\eta<\delta$,
\begin{align}\label{integral:bounds}
 \int_{\{U>\delta\}}e^{-U/\eps}\dd x
 \leq e^{-(\delta-\eta/2)/\eps}
       \int_{\R^d}e^{-\eta U/(2\delta\eps)}\dd x,
\end{align}
and
\begin{align*}
 Z_\eps\geq |A_\eta|e^{-\eta/(2\eps)}.
\end{align*}
The last integral in \eqref{integral:bounds} is uniformly bounded for all sufficiently
small $\eps$ by quadratic confinement and $U\geq0$. Dividing the two bounds
gives \eqref{eq:gibbstail}.
The first condition in \eqref{eq:coolingassumption} converts this exponential
bound into $(t+t_0)^{-(\delta-\eta)/E+o(1)}$. Passing to the limiting initial
law by Lemma~\ref{lem:closure} and taking $r_0$ arbitrarily close to $1-D/E$
proves \eqref{eq:mainrate}.
\end{proof}

\end{document}
